\chapter{Geometria rzutowa: punkty w nieskończoności, krzywe i przecięcia}
\label{ch:rzutowa}

\section{Po co domykać przestrzeń afiniczną?}

Na płaszczyźnie afinicznej $\A^2(k)$ proste $y=0$ i $y=1$ nie
mają wspólnego punktu. Z drugiej strony proste $y=0$ i $y=\varepsilon x$
przecinają się w początku układu współrzędnych dla każdego
$\varepsilon\ne0$. Wyniki dotyczące liczby punktów przecięcia
stają się przez to niewygodne: musimy osobno omawiać proste równoległe,
a dla krzywych wyższych stopni mogą znikać kolejne przecięcia.

W geometrii rzutowej dodajemy punkty reprezentujące \emph{kierunki}.
Nie należy przy tym wyobrażać sobie dodatkowej zwykłej prostej
w tej samej płaszczyźnie euklidesowej. Otrzymujemy nową przestrzeń,
w której punkt jest \emph{klasą niezerowego wektora}.

\begin{center}
\begin{tikzpicture}[>=Stealth,scale=.93]
  % Schematic of two parallel affine lines: equations agree with vertical coordinates.
  \draw[->,gray!65] (-3.2,-.55)--(-3.2,1.6) node[above] {$y$};
  \draw[->,gray!65] (-3.25,0)--(3.65,0) node[right] {$x$};
  \draw[blue!70!black,very thick] (-2.9,0)--(3.25,0);
  \draw[red!75!black,very thick] (-2.9,1)--(3.25,1);
  \node[blue!70!black,anchor=south west,fill=white,inner sep=1pt]
    at (-2.7,0.05) {$y=0$};
  \node[red!75!black,anchor=south west,fill=white,inner sep=1pt]
    at (-2.7,1.05) {$y=1$};
  \draw[->,dashed,gray!70] (3.3,0.02)--(4.35,.42);
  \draw[->,dashed,gray!70] (3.3,1)--(4.35,.63);
  \node[draw,circle,fill=white,inner sep=2pt] at (4.75,.54) {$p_\infty$};
\end{tikzpicture}

\small Rysunek 2.1. Szkic idei: dwie proste o tym samym kierunku mają
jeden wspólny punkt na prostej w nieskończoności. Nie jest to rysunek
w metryce przestrzeni rzutowej.
\end{center}

\section{Przestrzeń rzutowa jako zbiór kierunków}

\begin{definicja}[Przestrzeń rzutowa]
\label{def:projective-space}
Niech $k$ będzie ciałem, a $n\ge0$. Na zbiorze
$k^{n+1}\setminus\{0\}$ definiujemy relację
\[
 (a_0,\ldots,a_n)\sim(b_0,\ldots,b_n)
 \quad\Longleftrightarrow\quad
 \exists\lambda\in k^\times\ \forall i\in\{0,\ldots,n\}:\ b_i=\lambda a_i.
\]
Jej klasy oznaczamy $[a_0:\cdots:a_n]$. \emph{Przestrzenią
rzutową wymiaru $n$ nad $k$} na poziomie punktów $k$-wymiernych
nazywamy $\PP^n(k)=(k^{n+1}\setminus\{0\})/\!\sim$.
Współrzędne $[a_0:\cdots:a_n]$ nazywamy \emph{jednorodnymi}.
\end{definicja}

W definicji występuje kwantyfikator $\lambda\ne0$. Nie możemy
utożsamiać wszystkich niezerowych wektorów ze sobą: dwa wektory
reprezentują ten sam punkt tylko wtedy, gdy leżą na tej samej
jednowymiarowej podprzestrzeni liniowej. Wektor zerowy nie reprezentuje
żadnego punktu. Dlatego $[1:2]=[2:4]$ nad $\R$, ale $[1:0]\ne[0:1]$.

\begin{stwierdzenie}[Dobrze określona relacja]
Dla każdego ciała $k$ i $n\ge0$ relacja z definicji
\ref{def:projective-space} jest relacją równoważności.
\end{stwierdzenie}
\begin{proof}
Zwrotność otrzymujemy dla $\lambda=1$. Jeśli $b=\lambda a$
przy $\lambda\ne0$, to $a=\lambda^{-1}b$, więc relacja jest
symetryczna. Jeżeli $b=\lambda a$ oraz $c=\mu b$, to
$c=(\mu\lambda)a$, a $\mu\lambda\ne0$, co dowodzi
przechodniości.
\end{proof}

\begin{przyklad}[Prosta rzutowa i jeden punkt uzupełniający]
W $\PP^1(k)$ każdy punkt ma co najmniej jedną niezerową
współrzędną. Jeśli $a_0\ne0$, otrzymujemy
$[a_0:a_1]=[1:a_1/a_0]$, a więc punkt odpowiada dokładnie
jednemu $t\in k$. Jedyny punkt, którego nie można tak zapisać,
to $[0:1]$. Stąd bijekcja zbiorów
$\PP^1(k)\cong k\sqcup\{\infty\}$.

Przy $k=\R$ nie jest to jeszcze twierdzenie o homeomorfizmie:
trzeba osobno określić, czy używamy topologii Zariskiego,
czy topologii otrzymanej z rzeczywistej przestrzeni wektorowej.
\end{przyklad}

\begin{przyklad}[Projektowe spotkanie dwóch prostych]
W $\PP^2(k)$ ustalmy odtąd kolejność współrzędnych
$[X:Y:Z]$, gdzie $Z$ jest współrzędną homogenizującą.
Równanie afiniczne $y=0$ przepisuje się jako $Y=0$,
a $y=1$ jako $Y-Z=0$. Wspólne rozwiązania spełniają
$Y=Z=0$, co daje dokładnie punkt $[1:0:0]$.
Jest on niedostępny w części $Z\ne0$ i reprezentuje kierunek
osi $x$.
\end{przyklad}

\section{Mapy afiniczne przestrzeni rzutowej}

\begin{definicja}[Standardowe mapy afiniczne]
Niech $k$ będzie ciałem, a $n\ge1$ liczbą całkowitą.
Dla $i\in\{0,\ldots,n\}$ niech
$U_i=\{[a_0:\cdots:a_n]\in\PP^n(k):a_i\ne0\}$.
Definiujemy $\varphi_i:U_i\to k^n$ przez podzielenie
wszystkich współrzędnych oprócz $a_i$ przez $a_i$,
z zachowaniem ich kolejności.
\end{definicja}

\begin{stwierdzenie}[Mapy pokrywają przestrzeń]
Dla każdego ciała $k$ oraz $n\ge1$ zachodzi
$\PP^n(k)=\bigcup_{i=0}^n U_i$; każde $\varphi_i$
jest bijekcją. Na $U_i\cap U_j$ przejścia współrzędnych
są opisane ilorazami współrzędnych z nieznikającym mianownikiem.
\end{stwierdzenie}
\begin{proof}
Punkt rzutowy ma niezerową współrzędną, więc należy do
przynajmniej jednego $U_i$. Ilorazy $a_j/a_i$ nie zależą
od reprezentanta, ponieważ przemnożenie wszystkich współrzędnych
przez $\lambda\ne0$ skraca się w liczniku i mianowniku.
Odwrotność $\varphi_i$ otrzymujemy, wstawiając $1$ na
i-tej pozycji. Oba złożenia są identycznościami, gdyż w jedną
stronę dzielimy przez $a_i$, a w drugą wybieramy reprezentanta
z tą współrzędną równą $1$. Na przecięciu dwóch map trzeba
wykonać te operacje kolejno, dlatego powstają ilorazy.
\end{proof}

\begin{przyklad}[Zmiana map na $\PP^1$]
Dla $[s:t]\in\PP^1(k)$ mamy $U_s=\{s\ne0\}$ i
$U_t=\{t\ne0\}$. Współrzędna na pierwszej mapie to
$u=t/s$, a na drugiej $v=s/t$. Na przecięciu $u\ne0$
i zachodzi $v=1/u$. Nie jest to wielomian na całej prostej,
ale jest funkcją regularną na otwartym zbiorze, na którym
mianownik nie znika. To powód, dla którego w geometrii
rzutowej nie wystarczą globalne wielomiany afiniczne.
\end{przyklad}

Dla $[X:Y:Z]\in\PP^2(k)$ utożsamiamy
$U_Z=\{Z\ne0\}$ z $\A^2(k)$ przez
$[X:Y:Z]\mapsto(X/Z,Y/Z)$. Dopełnieniem tej części
jest \emph{prosta w nieskończoności} $H_\infty=\{Z=0\}$,
izomorficzna z $\PP^1(k)$ przez $[X:Y:0]\mapsto[X:Y]$.

\section{Wielomiany jednorodne i algebraiczne podzbiory rzutowe}

\begin{definicja}[Wielomian jednorodny]
Niech $k$ będzie ciałem, $n\ge0$ liczbą całkowitą
i $S=k[X_0,\ldots,X_n]$.
Wielomian $F\in S$ nazywamy \emph{jednorodnym stopnia $d$},
gdy jest sumą jednomianów o łącznym stopniu $d$; wielomian
zerowy traktujemy jako jednorodny dowolnego stopnia.
Dla niezerowego $F$ oznaczamy jego stopień przez $\deg F=d$.
\end{definicja}

\begin{lemat}[Dlaczego jednorodność jest konieczna]
Niech $k$ będzie ciałem, a $n,d\ge0$ liczbami całkowitymi.
Jeśli $F\in k[X_0,\dots,X_n]$ jest jednorodny stopnia $d$,
to dla każdego $a\in k^{n+1}$ i $\lambda\in k$ zachodzi
$F(\lambda a)=\lambda^dF(a)$. W szczególności warunek
$F(a)=0$ nie zależy od wybranego niezerowego reprezentanta
punktu $[a]\in\PP^n(k)$.
\end{lemat}
\begin{proof}
Dla jednomianu $cX_0^{e_0}\cdots X_n^{e_n}$
o $e_0+\cdots+e_n=d$ podstawienie $\lambda a_i$
mnoży wartość przez $\lambda^{\sum e_i}=\lambda^d$.
Po zsumowaniu jednomianów otrzymujemy ten sam wzór dla $F$.
Gdy $\lambda\ne0$, mamy $\lambda^d\ne0$,
więc $F(\lambda a)=0$ wtedy i tylko wtedy, gdy $F(a)=0$.
\end{proof}

\begin{uwaga}[Wielomian niejednorodny nie określa równania rzutowego]
Nad $k=\R$ warunek $X+1=0$ nie jest dobrze określony dla
$[Z:X]\in\PP^1(\R)$: reprezentant $(1,-1)$ go spełnia,
a przemnożony przez $2$ reprezentant $(2,-2)$ już nie.
Poprawne równanie jednorodne ma postać $X+Z=0$.
\end{uwaga}

\begin{definicja}[Zbiór algebraiczny rzutowy]
Niech $k$ będzie ciałem, $n\ge0$, $S=k[X_0,\dots,X_n]$ i niech $T\subseteq S$ będzie
rodziną wielomianów jednorodnych. Definiujemy
$\V_+(T)=\{[a]\in\PP^n(k):F(a)=0\text{ dla każdego }F\in T\}$.
Podzbiór $X\subseteq\PP^n(k)$ nazywamy \emph{rzutowym
zbiorem algebraicznym}, gdy $X=\V_+(T)$ dla pewnej takiej rodziny.
Ideał $I\subseteq S$ jest \emph{jednorodny}, gdy jest generowany
przez wielomiany jednorodne.
\end{definicja}

\begin{twierdzenie}[Topologia Zariskiego na przestrzeni rzutowej]
\label{tw:projective-zariski}
Dla każdego ciała $k$ i liczby całkowitej $n\ge0$ zbiory algebraiczne rzutowe
są zbiorami domkniętymi pewnej topologii na $\PP^n(k)$.
Dla ideałów jednorodnych $I,J\subseteq S=k[X_0,\dots,X_n]$ zachodzą
$\V_+(I+J)=\V_+(I)\cap\V_+(J)$ oraz
$\V_+(IJ)=\V_+(I)\cup\V_+(J)$.
\end{twierdzenie}
\begin{proof}
Wielomian zerowy określa całą przestrzeń, a stała $1$
zbiór pusty. Dowolne przecięcie zbiorów zerowych odpowiada
połączeniu rodzin wielomianów. Dla sumy dwóch zbiorów
można powtórzyć argument z twierdzenia~\ref{tw:zariski-top}:
$F(a)G(a)=0$ w ciele wymusza $F(a)=0$ lub $G(a)=0$.
Mamy więc wymaganą stabilność na przecięcia dowolne i sumy
skończone. Równości dotyczące ideałów wynikają z tych samych
własności, ponieważ ideały sumy i iloczynu ideałów jednorodnych
są również jednorodne.
\end{proof}

\begin{uwaga}[Punkty a ideał nieistotny]
Ideał $\mathfrak m_+=(X_0,\dots,X_n)$ nazywa się
\emph{ideałem nieistotnym} (ang. \emph{irrelevant ideal}).
Mamy $\V_+(\mathfrak m_+)=\varnothing$: wektor reprezentujący
punkt rzutowy nie może mieć wszystkich współrzędnych zerowych.
W przeciwieństwie do przypadku afinicznego nie każdy właściwy
ideał jednorodny ma zatem niepusty zbiór punktów rzutowych.
\end{uwaga}

\begin{stwierdzenie}[Topologia na mapach]
Niech $k$ będzie ciałem, $n\ge1$ i $i\in\{0,\dots,n\}$.
Wtedy bijekcja $\varphi_i:U_i\to k^n$
przenosi topologię Zariskiego indukowaną z $\PP^n(k)$
na topologię Zariskiego na $k^n$.
\end{stwierdzenie}
\begin{proof}
Wystarczy omówić $i=0$. Ograniczenie jednorodnego równania
$F(X_0,\dots,X_n)=0$ do $X_0\ne0$ jest równoważne
afinicznemu równaniu $f(x_1,\dots,x_n)=F(1,x_1,\dots,x_n)=0$.
Zatem przecięcie rzutowego zbioru algebraicznego z $U_0$
przechodzi w zbiór afinicznie domknięty. Odwrotnie,
dla wielomianu zerowego wybieramy jednorodny wielomian
zerowy. Dla niezerowego $f\in k[x_1,\dots,x_n]$ stopnia $d$
wielomian $X_0^df(X_1/X_0,\dots,X_n/X_0)$
(po rozwinięciu) jest jednorodny i wyznacza na $U_0$
taki sam zbiór zerowy jak $f$. Dowodzi to równości
topologii na podprzestrzeniach.
\end{proof}

\section{Homogenizacja i domknięcie rzutowe}

W tej części współrzędna homogenizująca $Z$ zajmuje
\emph{ostatnią} pozycję: punkty zapisujemy jako
$[X_1:\cdots:X_n:Z]$, a część afiniczna odpowiada $Z\ne0$.
W szczególności na płaszczyźnie mamy uzgodniony zapis $[X:Y:Z]$.

\begin{definicja}[Homogenizacja]
Niech $k$ będzie ciałem, a $n\ge1$ liczbą całkowitą.
Dla niezerowego wielomianu $f\in k[x_1,\ldots,x_n]$
stopnia całkowitego $d$ definiujemy
\[
 f^h(X_1,\ldots,X_n,Z)
 =Z^d f(X_1/Z,\ldots,X_n/Z),
\]
rozumiejąc prawą stronę po skróceniu jako wielomian;
$f^h$ jest jednorodny stopnia $d$. Dla $f=0$ przyjmujemy
$f^h=0$. Operację $F\mapsto F(x_1,\ldots,x_n,1)$
nazywamy \emph{dehomogenizacją} względem $Z$.
\end{definicja}

\begin{przyklad}[Dwa równoległe równania]
Na płaszczyźnie afinicznej $y=0$ i $y=1$ mają równania
$f=y$ i $g=y-1$. Ich homogenizacje względem $Z$ to
$F=Y$ oraz $G=Y-Z$.
W $\PP^2(k)$ przecinają się w punkcie $[1:0:0]$.
Zatem dodanie punktów w nieskończoności odzyskuje
przecięcie, którego brakowało na płaszczyźnie afinicznej.
\end{przyklad}

\begin{definicja}[Homogenizacja ideału]
Niech $k$ będzie ciałem, $n\ge1$, a
$I\subseteq k[x_1,\ldots,x_n]$ ideałem.
Przez $I^h$ oznaczamy ideał jednorodny w
$k[X_1,\ldots,X_n,Z]$ generowany przez wszystkie $f^h$
z $f\in I$ (a nie tylko przez homogenizacje dowolnie
wybranego układu generatorów $I$).
\end{definicja}

\begin{twierdzenie}[Domknięcie rzutowe zbioru afinicznego]
\label{tw:projective-closure}
Niech $k$ będzie algebraicznie domkniętym ciałem,
$n\ge1$, $X\subseteq\A^n(k)$ zbiorem algebraicznym,
a $\iota:\A^n(k)\to\PP^n(k)$ niech będzie odwzorowaniem
$\iota(a_1,\dots,a_n)=[a_1:\dots:a_n:1]$.
Wówczas domknięcie $\iota(X)$ w topologii Zariskiego jest równe
$\V_+(\I_k(X)^h)$.
\end{twierdzenie}
\begin{proof}
Niech $Y=\overline{\iota(X)}$.
Ponieważ $\V_+(\I_k(X)^h)$ jest domknięty i każdy
$f\in\I_k(X)$ spełnia $f^h(a,1)=f(a)=0$ dla $a\in X$,
otrzymujemy $Y\subseteq\V_+(\I_k(X)^h)$.

Dla inkluzji odwrotnej weźmy dowolny zbiór domknięty
$C=\V_+(T)$ zawierający $\iota(X)$.
Każdy jednorodny $F\in T$ stopnia $d$ ma dehomogenizację
$f(x)=F(x,1)$, która znika na $X$; zatem $f\in\I_k(X)$.
Jeśli $f\ne0$ ma stopień $m\le d$, to
$F=Z^{d-m}f^h$. Gdy $f=0$, jednorodność $F$ wymusza
$F=0$: w jednorodnym $F$ wszystkie jednomiany po podstawieniu
$Z=1$ są różnymi jednomianami w pozostałych zmiennych.
Stąd każdy $F\in T$ znika na $\V_+(\I_k(X)^h)$.
Mamy więc $\V_+(\I_k(X)^h)\subseteq C$.
Po przecięciu wszystkich domkniętych $C$ zawierających
$\iota(X)$ dostajemy $\V_+(\I_k(X)^h)\subseteq Y$.
\end{proof}

\begin{uwaga}[Czy naprawdę potrzeba tu algebraicznej domkniętości?]
Powyższy dowód samej tożsamości domknięcia działa dla
każdego ciała $k$, jeśli $\I_k(X)$ oznacza rzeczywisty
ideał znikania na punktach $k$-wymiernych. Założenie
algebraicznej domkniętości jest przydatne, gdy chcemy
zastąpić $\I_k(X)$ pierwiastnikiem pierwotnego ideału równań
na mocy Nullstellensatz. Nie należy tych dwóch stwierdzeń
utożsamiać.
\end{uwaga}

\begin{przyklad}[Pułapka: homogenizacja samych generatorów]
W $k[x,y]$ rozważmy ideał
$I=(x,xy-y)=(x,y)$. Równość wynika z
$y=xy-(xy-y)\in I$. Jego zbiór zerowy to tylko punkt $(0,0)$,
którego domknięcie rzutowe powinno być punktem $[0:0:1]$.

Jeśli jednak homogenizujemy \emph{tylko} podane dwa generatory,
otrzymamy w $k[X,Y,Z]$ ideał
$J=(X,XY-YZ)$. Z jego równań wynika $X=0$ i $YZ=0$,
a więc $\V_+(J)=\{[0:0:1],[0:1:0]\}$.
Drugi punkt jest sztucznym punktem w nieskończoności.
Homogenizacja także elementu $y\in I$ daje $Y$ i usuwa
ten niepożądany punkt. Właściwy ideał $I^h$ to $(X,Y)$.
\end{przyklad}

\begin{przyklad}[Parabola i jej punkt w nieskończoności]
Równanie $y=x^2$ homogenizuje się do $YZ=X^2$.
W części $Z=1$ odzyskujemy zwykłą parabolę, a dla $Z=0$
mamy $X=0$, więc jedyny punkt w nieskończoności to $[0:1:0]$.
Dla hiperboli $xy=1$ otrzymujemy $XY=Z^2$;
na prostej $Z=0$ mamy $XY=0$, czyli dwa punkty
$[0:1:0]$ i $[1:0:0]$.
\end{przyklad}

\section{Krzywe rzutowe, proste, stożkowe i parametryzacje}

\begin{definicja}[Krzywa płaska i jej stopień]
Niech $k$ będzie algebraicznie domkniętym ciałem.
Niezerowy jednorodny $F\in k[X,Y,Z]$ dodatniego stopnia $d$
wyznacza rzutową krzywą płaską $C=\V_+(F)\subseteq\PP^2(k)$.
Jeśli $F$ jest bezkwadratowy (żaden nierozkładalny czynnik
nie występuje z wykładnikiem większym od $1$), to $d$
nazywamy \emph{stopniem} takiej krzywej z jej naturalnym
zredukowanym równaniem. Jeśli $F$ ma powtórzone czynniki,
stopień $F$ uwzględnia również te krotności.
\end{definicja}

\begin{uwaga}[Dlaczego dodajemy warunek bezkwadratowości?]
Wielomiany $X$ i $X^3$ opisują ten sam \emph{zbiór} prostej,
ale ich stopnie są różne. Stopień samego zredukowanego
zbioru algebraicznego nie może więc zależeć od arbitralnie
wybranego równania z potęgami jego czynników.
\end{uwaga}

\begin{przyklad}[Stożkowa jest obrazem prostej rzutowej]
Niech $k$ będzie ciałem, a $C=\V_+(XZ-Y^2)\subseteq\PP^2(k)$.
Odwzorowanie
$\nu:\PP^1(k)\to C$, $[s:t]\mapsto[s^2:st:t^2]$
jest dobrze określone: po zamianie $(s,t)$ na
$(\lambda s,\lambda t)$ wszystkie trzy współrzędne
mnożą się przez $\lambda^2$. Równanie krzywej zachodzi,
bo $s^2t^2-(st)^2=0$.

Jeżeli $X\ne0$, to $s\ne0$, a parametr odtwarzamy przez
$[s:t]=[1:Y/X]$. Jeżeli $Z\ne0$, analogicznie
$[s:t]=[Y/Z:1]$. Na krzywej nie ma punktu z $X=Z=0$:
wtedy $Y^2=0$ i wszystkie współrzędne byłyby zerowe.
Formuły odwrotne mają regularne współrzędne na odpowiednich
mapach, zatem $C$ jest izomorficzna z $\PP^1$.
\end{przyklad}

\begin{center}
\begin{tikzpicture}[>=Stealth,scale=.83]
 \draw[->,gray!70] (-3.1,0)--(3.45,0) node[right] {$x$};
 \draw[->,gray!70] (0,-.5)--(0,4.5) node[above] {$y$};
 \draw[blue!70!black,thick,domain=-1.9:1.9,samples=100] plot (\x,{\x*\x});
 \fill[red!80!black] (0,0) circle(2.1pt);
 \draw[densely dashed,->] (2.15,4.05)--(3,4.4);
 \node[font=\small,align=left,anchor=west,fill=white,inner sep=2pt] at (3.5,3.3)
 {Jeden punkt\\w nieskończoności\\$[0:1:0]$};
 \node[blue!70!black] at (-2.05,3.9) {$YZ=X^2$};
\end{tikzpicture}

\small Rysunek 2.2. Afiniczny fragment stożkowej $YZ=X^2$ i kierunek
jej jedynego punktu w nieskończoności.
\end{center}

\section{Odwzorowania rzutowe i ich miejsca nieokreśloności}

\begin{definicja}[Odwzorowanie zadane formami]
Niech $k$ będzie ciałem, $n,m\ge0$, a
$F_0,\dots,F_m\in k[X_0,\dots,X_n]$ formami jednorodnymi tego samego dodatniego stopnia $d$.
Na zbiorze $U=\PP^n(k)\setminus\V_+(F_0,\dots,F_m)$
definiujemy odwzorowanie
$\Phi:U\to\PP^m(k)$ wzorem
$\Phi([a])=[F_0(a):\cdots:F_m(a)]$.
Zbiór $B=\V_+(F_0,\dots,F_m)$ nazywamy
\emph{zbiorem bazowym} tego zapisu.
\end{definicja}

\begin{stwierdzenie}[Niezależność od reprezentanta]
Niech $k$ będzie ciałem, $n,m\ge0$, a
$F_0,\ldots,F_m\in k[X_0,\ldots,X_n]$
formami jednorodnymi tego samego dodatniego stopnia $d$.
Jeśli $U=\PP^n(k)\setminus\V_+(F_0,\ldots,F_m)$, to wzór
$\Phi([a])=[F_0(a):\cdots:F_m(a)]$ wyznacza dobrze określone
odwzorowanie $\Phi:U\to\PP^m(k)$.
\end{stwierdzenie}
\begin{proof}
Na $U$ wartości $F_0(a),\dots,F_m(a)$ nie są jednocześnie
zerowe, więc definiują punkt rzutowy. Zamiana $a$ na
$\lambda a$ mnoży każdą z tych wartości przez $\lambda^d$,
co nie zmienia klasy jednorodnej w $\PP^m$.
\end{proof}

\begin{przyklad}[Przekształcenie Cremony]
Niech $k$ będzie ciałem. Na $\PP^2(k)$ rozważmy
$\Phi([X:Y:Z])=[YZ:XZ:XY]$.
Zbiór bazowy składa się z trzech punktów
$[1:0:0]$, $[0:1:0]$, $[0:0:1]$, bo wszystkie trzy iloczyny
są zerowe dokładnie wtedy, gdy co najwyżej jedna współrzędna
jest niezerowa. Na otwartym zbiorze $XYZ\ne0$ mamy
\[
 (\Phi\circ\Phi)([X:Y:Z])
 =[X^2YZ:XY^2Z:XYZ^2]=[X:Y:Z].
\]
Zatem $\Phi$ jest \emph{biwymiernym} przekształceniem
płaszczyzny rzutowej, ale nie jest zdefiniowana w trzech
punktach bazowych. Wymierne przekształcenie i wszędzie
określony morfizm nie są tym samym.
\end{przyklad}

\section{Punkty gładkie, styczne i osobliwości krzywych}

\begin{definicja}[Punkt gładki krzywej płaskiej]
Niech $k$ będzie ciałem algebraicznie domkniętym,
$F\in k[X,Y,Z]$ niezerowym jednorodnym wielomianem,
a $P=[a:b:c]\in\V_+(F)$. Punkt $P$ nazywamy
\emph{gładkim dla równania $F$}, jeśli przynajmniej jedna
z liczb $F_X(a,b,c)$, $F_Y(a,b,c)$, $F_Z(a,b,c)$ jest
niezerowa. W przeciwnym razie jest \emph{osobliwy dla $F$}.
Dla zredukowanej krzywej podane kryterium jest niezależne
od przemnożenia jej równania przez niezerową stałą.
\end{definicja}

\begin{uwaga}[Zredukowane równanie ma znaczenie]
Dla $F=X^2$ wszystkie pochodne zerują się na prostej $X=0$,
choć sama prosta zredukowana jest gładka.
Dlatego w dyskusji gładkości \emph{zbioru} krzywej wybieramy
jej zredukowane równanie, a nie dowolną potęgę.
\end{uwaga}

\begin{lemat}[Tożsamość Eulera]
Dla każdego ciała $k$, liczby całkowitej $n\ge0$
i jednorodnego wielomianu
$F\in k[X_0,\dots,X_n]$ stopnia $d$ zachodzi
$\sum_{i=0}^nX_i\frac{\partial F}{\partial X_i}=dF$.
\end{lemat}
\begin{proof}
Dla jednomianu $M=cX_0^{e_0}\cdots X_n^{e_n}$
wyrażenie po lewej jest równe
$(e_0+\cdots+e_n)M=dM$.
Równość zachowuje się przy sumowaniu, więc dotyczy $F$.
\end{proof}

\begin{twierdzenie}[Styczna rzutowa w punkcie gładkim]
\label{tw:tangent-projective}
Niech $k$ będzie ciałem algebraicznie domkniętym,
$F\in k[X,Y,Z]$ niezerową jednorodną formą,
a $P=[a:b:c]\in\V_+(F)$ punktem gładkim dla $F$.
Wtedy rzutowa prosta styczna do krzywej w $P$ ma równanie
\[
 F_X(a,b,c)X+F_Y(a,b,c)Y+F_Z(a,b,c)Z=0.
\]
Równanie to jest niezależne od reprezentanta punktu $P$.
\end{twierdzenie}
\begin{proof}
Co najmniej jeden współczynnik nie znika, więc wzór
opisuje prostą rzutową. Jeśli $F$ ma stopień $d$,
to zmiana reprezentanta $(a,b,c)$ na
$(\lambda a,\lambda b,\lambda c)$ mnoży pochodne
przez $\lambda^{d-1}$. Prosta się nie zmienia.
Tożsamość Eulera i $F(a,b,c)=0$ pokazują,
że punkt $P$ leży na podanej prostej.

Aby sprawdzić styczność, wybierzmy mapę, na której jedna
współrzędna $P$ jest niezerowa, i sprowadźmy ją do $1$.
W tej mapie równanie krzywej ma postać $f(u,v)=0$.
Wyraz liniowy rozwinięcia $f$ w punkcie $p=(u_0,v_0)$
wynosi $f_u(p)(u-u_0)+f_v(p)(v-v_0)$.
Jest to równanie afinicznej stycznej w gładkim punkcie.
Po przepisaniu do współrzędnych jednorodnych uzyskujemy
powyższy wzór; brakujący wyraz stały jest dokładnie
zapewniony przez tożsamość Eulera. W charakterystyce
dzielącej $d$ tożsamość Eulera wciąż obowiązuje,
choć $d=0$ jako element ciała, a rachunek lokalny
nadal daje poprawne równanie stycznej.
\end{proof}

\begin{przyklad}[Styczna do stożkowej]
Niech $k$ będzie ciałem. Dla $F=XZ-Y^2$ otrzymujemy $F_X=Z$, $F_Y=-2Y$,
$F_Z=X$. W punkcie $P=[1:t:t^2]$ prosta styczna ma
równanie $t^2X-2tY+Z=0$. Podstawienie tego punktu daje
$t^2-2t^2+t^2=0$, jak powinno. Gdy charakterystyka
wynosi $2$, współczynnik $-2t$ jest zerem i wzór
wciąż poprawnie opisuje styczną.
\end{przyklad}

\begin{przyklad}[Ostrze i węzeł]
Dla równania afinicznego $y^2=x^3$ obie pochodne
$-3x^2$ i $2y$ znikają w $(0,0)$; zatem początek
jest punktem osobliwym. Dla równania
$y^2=x^2(x+1)$ w początku również znikają obie pochodne.
Nad $\R$ pierwsza krzywa ma ostrze, a druga lokalnie
dwie gałęzie o różnych stycznych. To różne typy osobliwości,
choć w obu przypadkach kryterium Jacobiego wykrywa
niegładkość. Bardziej systematyczna teoria będzie
w następnym rozdziale o punktach osobliwych.
\end{przyklad}

\section{Przecięcia z prostą i pierwsza wersja twierdzenia Bézouta}

\begin{definicja}[Krotność pierwiastka formy binarnej]
Niech $k$ będzie algebraicznie domkniętym ciałem,
a $H(s,t)\in k[s,t]$ niezerową formą jednorodną stopnia $d$.
Każdy punkt $P\in\PP^1(k)$ odpowiada formie liniowej
$\ell_P(s,t)$, której zbiór zerowy w $\PP^1(k)$ jest
równy $\{P\}$. \emph{Krotnością zera $H$ w $P$} nazywamy
największe $m\ge0$ takie, że $\ell_P^m$ dzieli $H$.
\end{definicja}

\begin{lemat}[Rozkład formy binarnej]
Niech $k$ będzie ciałem algebraicznie domkniętym,
a $H\in k[s,t]$ niezerową formą jednorodną stopnia $d\ge1$.
Wtedy istnieją formy liniowe $\ell_1,\ldots,\ell_d$
i $c\in k^\times$ takie, że $H=c\ell_1\cdots\ell_d$.
W szczególności suma krotności wszystkich zer $H$ w $\PP^1(k)$
wynosi $d$.
\end{lemat}
\begin{proof}
Podstawmy $s=1$. Otrzymujemy wielomian
$h(t)=H(1,t)$ stopnia $m\le d$. Jest on niezerowy:
jednorodna forma $H$ jest jednoznacznie wyznaczona przez
współczynniki jednomianów $s^{d-j}t^j$, które po
podstawieniu $s=1$ stają się współczynnikami $t^j$.
Dzięki algebraicznej domkniętości mamy
$h(t)=c\prod_{j=1}^m(t-\alpha_j)$, z uwzględnieniem krotności.
Homogenizując z powrotem do stopnia $d$, mamy
$H(s,t)=c\,s^{d-m}\prod_{j=1}^m(t-\alpha_js)$.
Współczynnik przy $s^{d-m}$ opisuje zero $[0:1]$
o krotności $d-m$. Gdy $h$ jest stały, otrzymujemy
$H=cs^d$; przypadek $d=0$ nie jest tu potrzebny.
Suma wykładników czynników liniowych jest zawsze równa $d$.
\end{proof}

\begin{twierdzenie}[Bézout dla prostej i krzywej]
\label{tw:bezout-line}
Niech $k$ będzie algebraicznie domkniętym ciałem,
$F\in k[X,Y,Z]$ niezerową formą stopnia $d\ge1$,
a $L\subseteq\PP^2(k)$ prostą rzutową niebędącą
składową krzywej $\V_+(F)$. Wtedy $L\cap\V_+(F)$
jest niepusty i skończony, a suma krotności przecięć
w jego punktach wynosi $d$. Krotność w $P\in L$
definiujemy jako krotność zera otrzymanej formy binarnej
po parametryzacji liniowej $L\cong\PP^1(k)$.
\end{twierdzenie}
\begin{proof}
Prostą $L$ możemy sparametryzować jako
$[s:t]\mapsto[s\,u+t\,v]$, gdzie $u,v\in k^3$
są liniowo niezależne i rozpinają dwuwymiarową
podprzestrzeń odpowiadającą $L$. Po podstawieniu
otrzymujemy jednorodną formę $H(s,t)=F(su+tv)$ stopnia $d$.
Nie jest ona zerowa, ponieważ w przeciwnym razie $F$
znikałby na całej prostej $L$, a jej równanie liniowe
dzieliłoby $F$ (po liniowej zmianie współrzędnych
byłaby to podzielność przez jedną zmienną). To oznaczałoby,
że $L$ jest składową $\V_+(F)$, wbrew założeniu.
Na mocy poprzedniego lematu suma krotności zer $H$
wynosi $d$. Punkty zerowe $H$ odpowiadają dokładnie
punktom przecięcia $L$ i krzywej.
\end{proof}

\begin{przyklad}[Dlaczego styczność liczy się podwójnie?]
Dla stożkowej $XZ-Y^2=0$ i prostej $Z=X$ podstawiamy
$[X:Y:Z]=[s:t:s]$. Otrzymujemy $s^2-t^2$.
Jeśli $\operatorname{char}k\ne2$, są to dwa przecięcia
$[1:1:1]$ i $[1:-1:1]$, każde krotności $1$.
Natomiast prosta styczna w punkcie $[1:0:0]$ ma równanie
$Z=0$. Jej ograniczenie daje $-Y^2$, więc jedyny punkt
przecięcia $[1:0:0]$ ma krotność $2$.
\end{przyklad}

\section{Ogólne twierdzenie Bézouta: po co nam algebra ilorazowa?}

Dla dwóch krzywych stopni $d$ i $e$ oczekiwalibyśmy
$d\cdot e$ punktów przecięcia. Bez krotności jest to
fałszywe nawet dla prostej stycznej do stożkowej.
Potrzebujemy więc definicji uwzględniającej zarówno
zwykłe przecięcia, jak i styczność oraz punkty osobliwe.
Poniżej dowodzimy pełnej wersji płaskiego twierdzenia Bézouta
bez odwoływania się do nieudowodnionego wzoru na wymiar krzywej.
Przyda się jednak kilka lematów algebraicznych.

\begin{definicja}[Pierścień stopniowany i szereg Hilberta]
Niech $k$ będzie ciałem. Pierścień $S=k[X,Y,Z]$ ma \emph{stopniowanie}
$S=\bigoplus_{m\ge0}S_m$, gdzie $S_m$ składa się z form
jednorodnych stopnia $m$. Dla jednorodnego ideału $J$
iloraz $S/J$ dziedziczy stopniowanie.
Dla stopniowanej $k$-przestrzeni $M=\bigoplus_{m\ge0}M_m$
o skończenie wymiarowych składowych \emph{szeregiem Hilberta}
nazywamy formalny szereg
$H_M(t)=\sum_{m\ge0}(\dim_kM_m)t^m$.
\end{definicja}

\begin{stwierdzenie}[Szereg Hilberta przekroju dwóch form]
\label{st:hilbert-complete-intersection}
Niech $k$ będzie ciałem i niech $F,G\in S=k[X,Y,Z]$
będą niezerowymi formami odpowiednio stopni $d,e\ge1$,
bez wspólnego nierozkładalnego czynnika. Wtedy dla
$M=S/(F,G)$ zachodzi tożsamość szeregów formalnych
\[
 H_M(t)=\frac{(1-t^d)(1-t^e)}{(1-t)^3}.
\]
\end{stwierdzenie}
\begin{proof}
Ponieważ monomy $X^aY^bZ^c$, $a+b+c=m$, stanowią bazę
$S_m$, mamy $\dim_kS_m=\binom{m+2}{2}$, a więc
$H_S(t)=(1-t)^{-3}$ (to wynika z trzykrotnego
pomnożenia szeregu geometrycznego $\sum_{m\ge0}t^m$).
Mnożenie przez niezerowe $F$ jest różnowartościowe,
bo $S$ jest dziedziną. Z dokładnego ciągu przestrzeni
stopniowanych $0\to S(-d)\xrightarrow{\cdot F}S\to S/(F)\to0$
otrzymujemy $H_{S/(F)}(t)=(1-t^d)H_S(t)$.

Mnożenie przez $G$ na $S/(F)$ jest także różnowartościowe:
jeśli $Gh\in(F)$, to $F\mid Gh$. Wielomiany w $S$
mają jednoznaczny rozkład na czynniki nierozkładalne
(twierdzenie~\ref{tw:polynomial-UFD}),
a $F$ i $G$ nie mają wspólnego czynnika. Każdy czynnik
$F$ wraz ze swoją krotnością musi zatem dzielić $h$;
czyli $F\mid h$. Wobec tego $[h]=0$ w $S/(F)$.
Drugi dokładny ciąg,
$0\to(S/(F))(-e)\xrightarrow{\cdot G}S/(F)\to M\to0$,
daje $H_M(t)=(1-t^e)H_{S/(F)}(t)$, co kończy dowód.
\end{proof}

\begin{uwaga}[Dokładność i przesunięcie stopnia]
Przesunięcie $N(-d)$ oznacza ten sam moduł z nowym
stopniowaniem $(N(-d))_m=N_{m-d}$ (i $N_j=0$ dla $j<0$).
Ciąg $0\to A\to B\to C\to0$ jest \emph{dokładny},
gdy obraz każdego odwzorowania jest jądrem następnego.
Dla każdej składowej stopnia $m$ mamy wówczas
$\dim B_m=\dim A_m+\dim C_m$, stąd dodawanie
szeregów Hilberta. Nie używamy tu jeszcze teorii
homologii ani funktorów pochodnych.
\end{uwaga}

\begin{lemat}[Skończoność przecięcia i wybór prostej w nieskończoności]
\label{lem:finiteness-plane-intersection}
Niech $k$ będzie ciałem algebraicznie domkniętym, a niezerowe
formy $F,G\in k[X,Y,Z]$ nie mają wspólnego czynnika.
Wtedy $\V_+(F,G)\subseteq\PP^2(k)$ jest zbiorem skończonym.
Istnieje więc prosta rzutowa $L$ nieprzechodząca przez
żaden z tych punktów.
\end{lemat}
\begin{proof}
Wybierzmy punkt $Q\in\PP^2(k)$, który nie należy do
$\V_+(F)\cup\V_+(G)$. Taki punkt istnieje:
oba wielomiany są niezerowe, a nad nieskończonym ciałem
iloczyn $FG$ nie może znikać na całym $k^3$.
Po liniowej zmianie współrzędnych możemy przyjąć
$Q=[0:1:0]$, a wtedy współczynniki przy $Y^d$ w $F$
i przy $Y^e$ w $G$ są niezerowymi stałymi.
Postrzegamy $F,G$ jako wielomiany w $Y$ o współczynnikach
w $k[X,Z]$ i przechodzimy do ciała $K=k(X,Z)$.

Nie mają wspólnego czynnika w $K[Y]$. Rzeczywiście,
każdy z nich jest \emph{pierwotny} w $k[X,Z][Y]$,
ponieważ jego współczynnik wiodący jest niezerową stałą.
Gdyby miały wspólny niejednostkowy czynnik nad $K$,
po usunięciu mianowników i zastosowaniu lematu Gaussa
dla wielomianów nad pierścieniem z jednoznacznym rozkładem
otrzymalibyśmy wspólny czynnik w $k[X,Y,Z]$.

Dla dwóch wielomianów $A,B\in K[Y]$ stopni $d,e$
definiujemy \emph{rezultant} $\operatorname{Res}_Y(A,B)$
jako wyznacznik macierzy liniowego odwzorowania
\[
 K[Y]_{<e}\oplus K[Y]_{<d}\longrightarrow K[Y]_{<d+e},
 \qquad (U,V)\longmapsto UA+VB.
\]
Tutaj $K[Y]_{<r}$ oznacza wielomiany stopnia mniejszego niż $r$.
Jeśli wyznacznik jest zerowy, istnieje niezerowa para
$(U,V)$ z $UA+VB=0$. Dla wielomianów względnie pierwszych
wynikałoby $B\mid U$, co jest niemożliwe przy $\deg U<e$,
chyba że $U=0$, a wtedy także $V=0$; sprzeczność.
Odwrotnie, przy wspólnym czynniku $D$ można wybrać
$U=B/D$ i $V=-A/D$, więc wyznacznik jest zerowy.
Zatem rezultat $R(X,Z)=\operatorname{Res}_Y(F,G)$
jest \emph{niezerowym} wielomianem w $k[X,Z]$.

Jeśli $P=[x:y:z]$ jest wspólnym zerem $F,G$, to
$[x:z]\ne[0:0]$, gdyż inaczej $P=Q$, a $F(Q)G(Q)\ne0$.
Podstawienie $(X,Z)=(x,z)$ do macierzy rezultantu
daje dwa wielomiany w $Y$ o wspólnym pierwiastku $y$,
więc $R(x,z)=0$.
Wielomian $R$ jest jednorodny stopnia $de$.
Aby uzasadnić stopień, zapiszmy
$F=\sum_{j=0}^d f_j(X,Z)Y^j$ i
$G=\sum_{j=0}^e g_j(X,Z)Y^j$; współczynniki $f_j$
i $g_j$ są jednorodne stopni odpowiednio $d-j$ i $e-j$.
W macierzy rezultantu jest $e$ wierszy odpowiadających
$Y^iF$, $0\le i<e$, i $d$ wierszy odpowiadających
$Y^iG$, $0\le i<d$. W każdym składniku wyznacznika
wybiera się po jednym współczynniku z każdego wiersza
i z każdej kolumny potęg $Y^0,\dots,Y^{d+e-1}$.
Suma wykładników $Y$ po odjęciu przesunięć wierszy
wynosi zatem
$\sum_{j=0}^{d+e-1}j-\sum_{i=0}^{e-1}i-\sum_{i=0}^{d-1}i=de$.
Łączny stopień wybranych współczynników w $X,Z$
wynosi więc $ed+de-de=de$, niezależnie od składnika
wyznacznika.
Niezerowa forma binarna ma tylko skończenie wiele zer na
$\PP^1(k)$, jak wykazaliśmy wcześniej. Dla każdego z tych
skończenie wielu $[x:z]$ wielomian $F(x,Y,z)$ jest
niezerowy (ma niezerowy współczynnik przy $Y^d$),
więc ma skończenie wiele pierwiastków $y$.
W rezultacie $\V_+(F,G)$ jest skończony.

Ponieważ $k$ jest nieskończone, spośród prostych rzutowych
możemy wybrać prostą omijającą wszystkie te punkty:
każdy zadany punkt wyznacza tylko właściwy podzbiór
przestrzeni współczynników możliwych równań liniowych,
a skończona suma takich właściwych zbiorów nie wyczerpuje
$k^3$. Po zmianie współrzędnych prostą tę możemy przyjąć
jako $Z=0$.
\end{proof}

\begin{lemat}[Jednorodna algebra przecięcia ma rangę $de$]
\label{lem:bezout-rank}
Niech $k$ będzie ciałem algebraicznie domkniętym,
a $F,G\in k[X,Y,Z]$ będą formami stopni $d,e\ge1$
bez wspólnego czynnika. Załóżmy, że na prostej $Z=0$
nie ma ich wspólnego zera rzutowego.
Wtedy $A=k[x,y]/(f,g)$, gdzie
$f(x,y)=F(x,y,1)$, $g(x,y)=G(x,y,1)$,
jest skończenie wymiarową $k$-algebrą i
$\dim_kA=de$.
\end{lemat}
\begin{proof}
Niech $S=k[X,Y,Z]$, $M=S/(F,G)$ oraz
$\bar F=F(X,Y,0)$, $\bar G=G(X,Y,0)$.
Formy $\bar F,\bar G$ nie są zerowe: gdyby np.
$\bar F=0$, to $Z\mid F$. Ponieważ $F,G$ nie mają
wspólnego czynnika, koniecznie $Z\nmid G$,
czyli $\bar G\ne0$. Niezerowa forma $\bar G$
dodatniego stopnia ma nad ciałem algebraicznie
domkniętym zero na prostej $Z=0$, sprzecznie
z założeniem. Podobnie dla $\bar G=0$. Formy
$\bar F,\bar G$ są względnie pierwsze w $k[X,Y]$:
inaczej wspólny czynnik liniowy dawałby wspólny punkt
na prostej w nieskończoności.

Pokażemy, że mnożenie przez $Z$ jest różnowartościowe
w $M$. Jeśli $ZH=UF+VG$, to po podstawieniu $Z=0$
mamy $\bar U\bar F+\bar V\bar G=0$.
Ponieważ $\gcd(\bar F,\bar G)=1$, wynika stąd
$\bar G\mid\bar U$ i $\bar F\mid\bar V$.
Istnieje zatem $\bar W\in k[X,Y]$ z
$\bar U=\bar W\bar G$ i $\bar V=-\bar W\bar F$.
Wybierając dowolne podniesienie $W\in S$, możemy zapisać
$U=WG+ZU_1$, $V=-WF+ZV_1$. Otrzymujemy
$ZH=Z(U_1F+V_1G)$, a po skróceniu $Z$
w dziedzinie $S$ mamy $H\in(F,G)$.

Zatem mamy dokładny ciąg stopniowany
$0\to M(-1)\xrightarrow{\cdot Z}M\to M/ZM\to0$.
Szereg Hilberta z poprzedniego stwierdzenia daje
\[
 H_{M/ZM}(t)=(1-t)H_M(t)
 =\frac{(1-t^d)(1-t^e)}{(1-t)^2}
 =(1+t+\cdots+t^{d-1})(1+t+\cdots+t^{e-1}).
\]
Ponieważ $M/ZM\cong k[X,Y]/(\bar F,\bar G)$,
jest to skończenie wymiarowa $k$-przestrzeń, a jej
wymiar jest sumą współczynników podanego wielomianu,
czyli $de$.

Wyjaśnijmy, dlaczego stąd wynika $\dim_kM/(Z-1)M=de$.
Wybierzmy jednorodne elementy $b_1,\ldots,b_{de}\in M$
reprezentujące bazę stopniowanego ilorazu $M/ZM$.
Pokażemy, że są bazą $M$ jako modułu nad $k[Z]$.
Dla dowolnego elementu jednorodnego $h\in M$ jego klasa
modulo $ZM$ jest kombinacją klas tych $b_i$, których
stopień jest równy $\deg h$. Po odjęciu tej kombinacji
pozostaje $Zh_1$ z $\deg h_1=\deg h-1$.
Indukcja względem stopnia pokazuje, że $b_i$ generują $M$
nad $k[Z]$. Mamy więc surjekcję stopniowanych modułów
$\bigoplus_i k[Z](-\deg b_i)\twoheadrightarrow M$.
Szereg Hilberta źródła wynosi
$\sum_i t^{\deg b_i}/(1-t)=H_{M/ZM}(t)/(1-t)$,
czyli $H_M(t)$ dzięki dokładnemu ciągowi powyżej.
W każdej składowej stopnia mapa jest więc surjekcją
między przestrzeniami tego samego skończonego wymiaru;
jest izomorfizmem. Dlatego $M\cong\bigoplus_i k[Z]b_i$.

Iloraz przez $Z-1$ zamienia każdy składnik $k[Z]b_i$
w jednowymiarową przestrzeń nad $k$, a ponadto
$M/(Z-1)M\cong k[x,y]/(f,g)=A$, ponieważ podstawienie
$Z=1$ dehomogenizuje $F$ i $G$. Stąd $\dim_kA=de$.
\end{proof}

\begin{definicja}[Lokalna krotność przecięcia]
\label{def:intersection-multiplicity}
Niech $k$ będzie ciałem algebraicznie domkniętym,
a $f,g\in k[x,y]$ nie mają wspólnej składowej algebraicznej
i mają skończenie wiele wspólnych zer.
Załóżmy, że algebra $A=k[x,y]/(f,g)$ jest skończenie wymiarowa
nad $k$. Dla wspólnego zera $p=(a,b)$ oznaczmy przez
$\mathfrak m_p=(x-a,y-b)A$ ideał maksymalny $A$.
Przez $A_{\mathfrak m_p}$ oznaczamy lokalizację $A$
względem zbioru $A\setminus\mathfrak m_p$:
jej elementy to ułamki $u/s$ z $s\notin\mathfrak m_p$,
z relacją zwykłego mnożenia na krzyż po ewentualnym
pomnożeniu przez dodatkowy mianownik.
Definiujemy \emph{lokalną krotność przecięcia}
$I_p(f,g):=\dim_k A_{\mathfrak m_p}$.
\end{definicja}

\begin{lemat}[Algebra skończona rozkłada się na części lokalne]
\label{lem:artin-decomposition}
Niech $k$ będzie algebraicznie domkniętym ciałem,
a $A$ niezerową, skończenie wymiarową przemienną $k$-algebrą.
Wtedy $A$ ma skończenie wiele ideałów maksymalnych
$\mathfrak m_1,\ldots,\mathfrak m_r$ i
$\dim_kA=\sum_{i=1}^r\dim_k A_{\mathfrak m_i}$.
\end{lemat}
\begin{proof}
Ponieważ $A$ jest skończenie wymiarowa nad $k$, jest
pierścieniem noetherowskim. Jest też artinowski:
każdy malejący ciąg ideałów stabilizuje się, ponieważ
ideały są podprzestrzeniami $k$-liniowymi o nierosnących
wymiarach. Każdy jego iloraz będący dziedziną jest ciałem:
dla $0\ne a$ w skończenie wymiarowej dziedzinie
mnożenie przez $a$ jest różnowartościowym endomorfizmem
przestrzeni wektorowej, więc jest też surjektywne
i istnieje element odwrotny do $a$.
Zatem każdy ideał pierwszy $A$ jest maksymalny.

Każdy ideał maksymalny $\mathfrak m$ ma iloraz $A/\mathfrak m$
będący skończenie wymiarowym rozszerzeniem ciała $k$.
Algebraiczna domkniętość i lemat~\ref{lem:extension-trivial}
dają $A/\mathfrak m\cong k$.
Różne ideały maksymalne są parami współmaksymalne.
Dla dowolnych $r$ różnych ideałów maksymalnych chińskie
twierdzenie o resztach daje surjekcję $A\to k^r$,
a więc $r\le\dim_k A$. Liczba takich ideałów jest
zatem skończona.

Niech $N$ będzie nilradykałem $A$, czyli ideałem wszystkich
elementów nilpotentnych. Standardową równość
$N=\bigcap_i\mathfrak m_i$ uzasadnimy następująco:
każdy nilpotent leży w każdym ideale pierwszym;
odwrotnie, jeśli $a$ nie jest nilpotentny, zbiór
$\{1,a,a^2,\dots\}$ nie zawiera zera, a więc istnieje
ideał maksymalny w rodzinie ideałów rozłącznych z tym
zbiorem (lemat o maksymalności można otrzymać z
noetherowskości przez argument rosnącego łańcucha).
Taki maksymalny ideał rozłączny $J$ jest pierwszy.
Istotnie, jeśli $uv\in J$, lecz $u,v\notin J$, to oba
większe ideały $(J,u)$ i $(J,v)$ przecinają zbiór
$\{1,a,a^2,\dots\}$. Istnieją zatem liczby $r,s\ge0$
i elementy $j_1,j_2\in J$ oraz $b,c\in A$ takie, że
$a^r=j_1+bu$ i $a^s=j_2+cv$.
Po wymnożeniu otrzymujemy
$a^{r+s}=j_1j_2+j_1cv+j_2bu+bcuv\in J$,
co przeczy rozłączności $J$ ze zbiorem potęg $a$.
Ponieważ w $A$ wszystkie ideały pierwsze są maksymalne,
otrzymujemy wymaganą równość.

Ideał $N$ jest nilpotentny. Jeśli $N=(0)$, wystarcza
$N^1=(0)$. W przeciwnym przypadku jest skończenie
generowany przez niepusty układ $n_1,\dots,n_s$,
a każdy $n_i$ ma potęgę $n_i^{q_i}=0$.
W każdym iloczynie $\sum_i(q_i-1)+1$ generatorów
któryś $n_i$ wystąpi co najmniej $q_i$ razy,
więc iloczyn jest zerowy.
Stąd $N^M=0$ dla odpowiednio dużego $M$.
Ponieważ ideały $\mathfrak m_i$ są parami współmaksymalne,
$N=\bigcap_i\mathfrak m_i=\prod_i\mathfrak m_i$.
Wynika stąd $\prod_i\mathfrak m_i^M=0$.
Ideały $\mathfrak m_i^M$ też są parami współmaksymalne,
więc chińskie twierdzenie o resztach daje
$A\cong\prod_i A/\mathfrak m_i^M$.

W ilorazie $B_i=A/\mathfrak m_i^M$ jedynym ideałem maksymalnym
jest obraz $\mathfrak m_i$: każdy inny ideał maksymalny
$\mathfrak m_j$ nie zawiera $\mathfrak m_i^M$.
Ponadto każdy element spoza tego ideału jest odwracalny:
pozostaje poza wszystkimi ideałami maksymalnymi, a element
nieodwracalny zawsze należy do pewnego ideału maksymalnego.
Zatem $B_i$ jest pierścieniem lokalnym.
Lokalizacja iloczynu $\prod_j B_j$ względem dopełnienia
i-tego ideału maksymalnego usuwa wszystkie pozostałe czynniki
(idempotenty w nich stają się zerami), a zachowuje $B_i$.
Dlatego $A_{\mathfrak m_i}\cong B_i$. Biorąc wymiary
wektorowe rozkładu produktu, otrzymujemy tezę.
\end{proof}

\begin{twierdzenie}[Bézout dla dwóch krzywych płaskich]
\label{tw:bezout-general}
Niech $k$ będzie ciałem algebraicznie domkniętym,
a $F,G\in k[X,Y,Z]$ niezerowymi wielomianami jednorodnymi
stopni odpowiednio $d,e\ge1$, bez wspólnego
nierozkładalnego czynnika. Wtedy wspólnych punktów
krzywych $\V_+(F)$ i $\V_+(G)$ jest skończenie wiele,
a suma ich lokalnych krotności przecięcia wynosi $de$:
\[
 \sum_{P\in\V_+(F,G)} I_P(F,G)=de.
\]
Lokalną krotność definiujemy w dowolnej mapie afinicznej
zawierającej $P$ przez lokalizację ilorazu odpowiednich
dehomogenizacji, jak w definicji~\ref{def:intersection-multiplicity}.
\end{twierdzenie}
\begin{proof}
Z lematu~\ref{lem:finiteness-plane-intersection} zbiór
wspólnych punktów jest skończony i można wybrać prostą
w nieskończoności, która go omija. Po rzutowej liniowej
zmianie współrzędnych możemy przyjąć, że tą prostą
jest $Z=0$. Każdy wspólny punkt leży więc w mapie
$Z\ne0$ i odpowiada wspólnemu zeru wielomianów
$f(x,y)=F(x,y,1)$ i $g(x,y)=G(x,y,1)$.

Na mocy lematu~\ref{lem:bezout-rank} algebra
$A=k[x,y]/(f,g)$ ma wymiar $de$ nad $k$.
Jej ideały maksymalne odpowiadają dokładnie wspólnym
punktom $p=(a,b)$: wynika to z opisu maksymalnych ideałów
w $k[x,y]$ nad ciałem algebraicznie domkniętym
(wniosek~\ref{wn:maximal-point}) oraz z twierdzenia
o odpowiedniości ideałów w ilorazie. Wreszcie lemat
\ref{lem:artin-decomposition} daje
$de=\dim_k A=\sum_{p}\dim_k A_{\mathfrak m_p}
=\sum_p I_p(f,g)$. Przy zmianie mapy rzutowej
lokalne pierścienie ilorazowe są izomorficzne.
Istotnie, w przecięciu map $X_i\ne0$ i $X_j\ne0$
zmiana współrzędnych jest regularna i ma regularną odwrotność.
Jeśli $F$ ma stopień $d$, jego dehomogenizacje na tych
mapach spełniają $F/X_j^d=(X_i/X_j)^d(F/X_i^d)$;
analogicznie dla $G$ stopnia $e$. Czynniki $X_i/X_j$
i $X_j/X_i$ są na przecięciu odwracalne, więc
oba układy równań generują ten sam ideał po
przeniesieniu do wspólnego pierścienia funkcji
regularnych. Lokalizacja w punkcie $P$ daje zatem
izomorficzne lokalne pierścienie ilorazowe.
Ich wymiar nad $k$, a więc lokalna krotność przecięcia,
nie zależy od wybranej mapy. To dowodzi tezy.
\end{proof}

\begin{przyklad}[Dwa przecięcia, łącznie stopień cztery]
Niech $k$ będzie ciałem algebraicznie domkniętym
charakterystyki różnej od $2$. Dla $F=XZ-Y^2$
i $G=XZ+Y^2$ oba wielomiany mają
stopień $2$. Z równań $XZ-Y^2=XZ+Y^2=0$
otrzymujemy $Y^2=0$ oraz $XZ=0$, czyli punkty
$P=[1:0:0]$ i $Q=[0:0:1]$. W mapie $X=1$ wokół $P$
równania to $z-y^2=0$ i $z+y^2=0$.
W pierścieniu lokalnym iloraz jest izomorficzny
$k[y]_{(y)}/(y^2)$ (ponieważ $2$ jest odwracalne),
o wymiarze $2$ nad $k$. Analogicznie w mapie $Z=1$
wokół $Q$ krotność wynosi $2$. Suma $2+2=4$ zgadza się
z $\deg F\deg G=2\cdot2$.
\end{przyklad}

\section{Co należy zapamiętać przed teorią osobliwości i krzywych?}

Punkty przestrzeni rzutowej to kierunki w przestrzeni o jeden
wymiar większej. Równania geometryczne są jednorodne,
ponieważ tylko wtedy ich zbiory zerowe nie zależą od
reprezentanta punktu. Mapy $U_i\cong\A^n$ pozwalają
prowadzić zwykłe rachunki w lokalnych współrzędnych.
Przejście od zbioru afinicznego do jego domknięcia rzutowego
wymaga homogenizacji \emph{całego ideału znikania}.

Przecięcia rzutowe nie gubią punktów w nieskończoności,
ale ich liczba musi być liczona z krotnościami.
Dla prostej i krzywej stopnia $d$ dostajemy $d$,
a dla dwóch krzywych bez wspólnej składowej stopni $d,e$
dostajemy $de$. Ta sama algebra, która wcześniej służyła
do badania ideałów, pozwala teraz rachować przecięcia.
Kolejny naturalny temat to punkty osobliwe, stożki styczne,
normalizacja krzywych i ich rodzaj topologiczny.

% ======================================================================
% ROZDZIAŁ 3: PUNKTY OSOBLIWE I GEOMETRIA LOKALNA KRZYWYCH
% ======================================================================
